\chapter{PENDAHULUAN}

\section{Latar Belakang}

Model bahasa menghasilkan teks sedikit demi sedikit dalam satuan yang disebut token (potongan teks). Pada setiap langkah inferensi, yaitu saat model digunakan untuk menghasilkan keluaran, model menghitung representasi melalui lapisan-lapisan Transformer, sebuah arsitektur jaringan saraf, sebelum memilih token berikutnya. Penghitungan berulang ini dapat membuat proses inferensi memerlukan waktu dan komputasi yang besar. FiRST, misalnya, mengkaji pengurangan waktu inferensi dengan melewati sebagian lapisan berdasarkan masukan \citep{jain2025first}. Namun, pengurangan komputasi perlu dinilai bersama perubahan kualitas teks yang dihasilkan.

Salah satu cara mengurangi penghitungan tersebut adalah \textit{layer skipping}, yaitu melewati lapisan tertentu saat model menghasilkan keluaran. Pada strategi statis, lapisan yang dijalankan sama untuk semua masukan. Pada strategi adaptif, keputusan lapisan dapat bergantung pada isi masukan. PuDDing menggunakan router, yaitu komponen ringan yang memilih lapisan untuk dijalankan, berdasarkan prompt (teks masukan yang berisi instruksi dan bahan untuk diproses) \citep{wee2025pudding}. Hasil PuDDing pada tugas penalaran akal sehat yang diuji dalam penelitian tersebut menjadi bukti bagi pendekatan pemilihan adaptif pada tugas-tugas itu, tetapi tidak memastikan hasil yang sama pada peringkasan berita bahasa Indonesia.

Penelitian ini mengkaji pemilihan jalur eksekusi pada tingkat urutan masukan dan keluaran (\textit{sequence-level routing}). Router membaca representasi artikel satu kali untuk menentukan apakah setiap kelompok lapisan tengah dijalankan atau dilewati. Keputusan yang sama dipertahankan selama sisa pemrosesan prompt dan seluruh pembentukan token keluaran (\textit{decoding}), sehingga jalur tidak dipilih ulang pada setiap token. Pelatihan dilakukan dalam dua tahap, mengikuti pola yang digunakan FiRST \citep{jain2025first}. Pada tahap pertama, parameter model utama dibekukan dan hanya router yang dilatih. Pada tahap kedua, router dibekukan dan adapter Low-Rank Adaptation (LoRA) dilatih pada modul attention dan FFN untuk mengompensasi dampak pelompatan lapisan. Target router dibentuk dari ringkasan acuan dengan membandingkan kerugian prediksi token untuk sejumlah pola kelompok lapisan.

Model utama (\textit{backbone}) yang digunakan adalah Qwen3-0.6B, yaitu model bahasa kausal yang memprediksi token berikutnya berdasarkan teks sebelumnya, dengan ukuran nominal 0,6 miliar parameter dan 28 lapisan \citep{qwen2025modelcard}. Tugas penelitian adalah peringkasan berita bahasa Indonesia menggunakan bagian berbahasa Indonesia dari XL-Sum. Dataset XL-Sum berisi pasangan artikel dan ringkasan yang disiapkan untuk penelitian peringkasan abstraktif, yaitu peringkasan yang merumuskan kembali informasi penting dengan kalimat baru \citep{hasan2021xlsum}. Ringkasan acuan digunakan saat membentuk target pelatihan router, tetapi tidak diberikan kepada router sebagai masukan saat inferensi.

Jumlah lapisan yang dilewati saja belum menunjukkan manfaat praktis. Router menambah pekerjaan komputasi untuk membaca representasi prompt dan menentukan jalur. Karena itu, penelitian ini akan mengukur waktu inferensi aktual, termasuk biaya tambahan router, dan menilai perubahan kualitas ringkasan. Analisis tersebut diperlukan untuk melihat apakah penghematan waktu sebanding dengan perubahan kualitas.

Evaluasi membandingkan tiga pendekatan utama: model penuh, pelompatan lapisan statis, dan routing adaptif. LoRA dilatih pada masing-masing pendekatan; routing adaptif tanpa LoRA juga diuji sebagai kontrol untuk menilai kontribusi LoRA. Pola statis dipilih menggunakan data validasi, lalu diterapkan pada seluruh artikel. Efisiensi utama dibandingkan berdasarkan waktu inferensi aktual pada lingkungan dan pengaturan yang sama. Jumlah lapisan aktif dan panjang keluaran dilaporkan sebagai informasi tambahan. Penelitian ini tidak mengasumsikan bahwa routing atau LoRA akan lebih unggul. Hasil evaluasi digunakan untuk menganalisis perubahan kualitas dan biaya inferensi pada setiap kondisi.

\section{Rumusan Masalah}

Berdasarkan latar belakang tersebut, penelitian ini mengajukan pertanyaan berikut.

\begin{enumerate}
  \item Bagaimana pengaruh sequence-level routing terhadap kualitas ringkasan dan waktu inferensi Qwen3-0.6B pada XL-Sum Indonesia?
  \item Apakah routing adaptif dengan LoRA memberikan trade-off kualitas--efisiensi yang lebih baik daripada pelompatan lapisan statis dan model penuh dengan LoRA, serta bagaimana hasil routing adaptif berubah ketika LoRA tidak digunakan?
\end{enumerate}

\section{Batasan Masalah}

\subsection{Cakupan Penelitian}

Penelitian ini dibatasi pada aspek-aspek berikut.

\begin{enumerate}
  \item Model yang digunakan adalah Qwen3-0.6B, sebuah model bahasa kecil berbasis decoder (\textit{decoder-only small language model}, SLM).
  \item Tugas penelitian adalah peringkasan berita berbahasa Indonesia menggunakan subset Indonesia dari XL-Sum.
  \item Mekanisme yang ditambahkan terdiri atas router ringan dan pelaksanaan kelompok lapisan secara kondisional. Sebanyak 28 lapisan dibagi menjadi tujuh kelompok berurutan, masing-masing berisi empat lapisan. Kelompok pertama dan terakhir selalu dijalankan. Setiap kelompok tengah dapat dijalankan atau dilewati; jumlah lapisan aktif tidak ditetapkan sama untuk setiap artikel.
  \item Router menggunakan rata-rata vektor representasi internal (\textit{hidden state}) token artikel setelah kelompok pertama. Token instruksi, token khusus, dan padding tidak disertakan dalam rata-rata. Sebuah jaringan perceptron multilapis (\textit{multilayer perceptron}, MLP) ringan menghasilkan keputusan biner untuk masing-masing kelompok tengah. Rincian arsitektur dan konfigurasi router akan ditetapkan pada tahap implementasi dan pemilihan konfigurasi menggunakan data validasi. Jalur dipilih berdasarkan prompt dan dipertahankan selama sisa pemrosesan prompt serta seluruh pembentukan token keluaran.
  \item Pelatihan terdiri atas dua tahap. Pada tahap pertama, parameter Qwen dibekukan dan hanya parameter router yang diperbarui. Target router dibentuk pada data latih melalui pencarian bertahap yang menyeimbangkan kerugian prediksi token dan proporsi kelompok tengah aktif. Router kemudian dilatih untuk memprediksi target tersebut menggunakan fungsi kerugian klasifikasi biner. Pada tahap kedua, bobot router dibekukan dan adapter LoRA dilatih pada modul attention dan FFN di seluruh lapisan Qwen menggunakan data latih. Pada routing adaptif, objektif pelatihan LoRA juga dirancang untuk mempertahankan manfaat pelompatan lapisan. Bentuk operasional penalti dan konfigurasi LoRA yang belum ditetapkan akan dipilih atau diverifikasi menggunakan data validasi sebelum evaluasi akhir.
  \item Evaluasi mencakup tiga pendekatan utama: model penuh+LoRA, pelompatan statis+LoRA, dan routing adaptif+LoRA. Routing adaptif tanpa LoRA menjadi kontrol tambahan. Satu pola statis dipilih menggunakan data validasi dan diterapkan pada seluruh artikel. Kualitas dan efisiensi dibandingkan menggunakan protokol yang sama; waktu inferensi aktual menjadi ukuran efisiensi utama, sedangkan jumlah lapisan aktif dan panjang keluaran menjadi informasi tambahan.
  \item Kualitas ringkasan dinilai sepenuhnya secara otomatis menggunakan skor F1 pada ROUGE-1, ROUGE-2, ROUGE-L, dan BERTScore. ROUGE mengukur kecocokan kata dan urutannya; BERTScore membandingkan representasi kontekstual teks. Skor F1 menggabungkan presisi dan cakupan. F1 dihitung untuk setiap artikel, lalu dirata-ratakan secara terpisah untuk tiap metrik; keempat skor dilaporkan sendiri-sendiri. Tidak dilakukan penilaian manual atau pemeriksaan faktualitas manual pada sampel. Efisiensi dinilai melalui waktu total per artikel, waktu pemrosesan prompt, waktu pembentukan keluaran, dan jumlah token keluaran. Pengukuran menggunakan ukuran batch satu, pemanasan eksekusi (\textit{warm-up}), dan pengulangan. Semua metode menggunakan mode non-thinking, yaitu mode yang mencegah Qwen menghasilkan teks penalaran sebelum jawaban akhir, dan konfigurasi \textit{greedy decoding} yang sama: \texttt{do\_sample=False} dan \texttt{num\_beams=1}. Perangkat dan presisi numerik juga akan disamakan.
  \item Batas jendela konteks eksperimen adalah 4.096 token untuk masukan dan keluaran secara gabungan. Keluaran dibatasi hingga maksimum 256 token, sehingga masukan dapat memuat paling banyak 3.840 token, termasuk instruksi dan format chat. Generasi berhenti saat model mengeluarkan token akhir respons sesuai konfigurasi model atau mencapai 256 token. Token akhir respons termasuk dalam batas 256 token dan tidak ditetapkan panjang minimum. Batas dan aturan penghentian ini sama untuk seluruh metode.
\end{enumerate}

\subsection{Pengolahan dan Pemilihan Data}

Sebelum pemrosesan, spasi dan baris kosong dihapus hanya dari awal dan akhir artikel serta ringkasan acuan. Huruf besar, tanda baca, spasi di antara kata, dan isi teks lainnya dipertahankan. Pasangan artikel--ringkasan dikeluarkan jika salah satu teks kosong setelah pembersihan ini. Aturan tersebut berlaku pada data latih (\textit{train}), validasi (\textit{validation}), dan uji (\textit{test}); jumlah pasangan kosong yang dikeluarkan akan dicatat.

Artikel dianggap duplikat jika teksnya sama persis setelah pembersihan. Dalam satu bagian data (\textit{split}), kemunculan pertama menurut urutan asli dataset dipertahankan bersama ringkasan acuannya. Jika artikel yang sama muncul pada beberapa split, pasangan pada data uji diprioritaskan, diikuti data validasi lalu data latih. Salinan pada split lain dikeluarkan dan tidak dipindahkan. Identitas serta jumlah pasangan yang dikeluarkan akan dicatat.

Artikel dikeluarkan jika prompt lengkap yang memuat artikel utuh melebihi 3.840 token menurut tokenizer Qwen3-0.6B. Prompt mencakup instruksi dan format chat; artikel tidak dipotong. Ringkasan acuan juga dikeluarkan jika melebihi 256 token, termasuk token akhir respons menurut tokenizer yang sama. Ringkasan tidak dipotong. Kedua batas diterapkan pada seluruh split. Himpunan artikel evaluasi akan sama untuk semua metode. Jumlah artikel yang dikeluarkan karena batas panjang dan jumlah yang tersisa pada setiap split akan dicatat.

Jika hanya sebagian data digunakan, sampel dipilih secara acak tanpa pengembalian pada setiap split setelah pembersihan, penyaringan, dan penghapusan duplikasi. Jumlah artikel pada setiap split akan ditetapkan setelah waktu pemrosesan dan kebutuhan memori diperiksa. Pemilihan menggunakan seed tetap 42. Identitas sampel terpilih akan disimpan, dan semua metode akan menggunakan subset evaluasi yang sama.

\section{Tujuan Penelitian}

Penelitian ini bertujuan untuk:

\begin{enumerate}
  \item menganalisis pengaruh sequence-level routing terhadap kualitas ringkasan dan waktu inferensi Qwen3-0.6B pada XL-Sum Indonesia; dan
  \item membandingkan keseimbangan kualitas dan efisiensi (\textit{trade-off}) routing adaptif+LoRA dengan static skipping+LoRA dan model penuh+LoRA, serta mengukur kontribusi LoRA melalui kontrol router tanpa LoRA.
\end{enumerate}

\section{Manfaat Penelitian}

Secara akademik, penelitian ini diharapkan memberikan bukti empiris tentang perilaku routing adaptif pada Qwen3-0.6B untuk peringkasan berita bahasa Indonesia. Hasilnya akan menunjukkan hubungan antara kelompok lapisan yang dijalankan, kualitas ringkasan, dan waktu inferensi pada rancangan yang diuji. Temuan dan keterbatasannya dapat membantu merumuskan pertanyaan untuk penelitian lanjutan.

Secara praktis, dokumentasi implementasi dan protokol evaluasi diharapkan membantu peneliti lain menguji pelaksanaan lapisan secara kondisional pada model ini. Hasil penelitian dapat digunakan untuk menilai penghematan waktu inferensi beserta perubahan kualitas ringkasan, tanpa menganggap bahwa melewati lapisan selalu mempercepat proses.

\section{Sistematika Penulisan}

Proposal ini terdiri atas empat bab. Bab I menjelaskan latar belakang, rumusan masalah, batasan, tujuan, dan manfaat penelitian. Bab II mengulas penelitian terkait dan menjelaskan posisi penelitian ini. Bab III memaparkan metode, rancangan sistem, tahapan penelitian, dan rencana evaluasi. Bab IV menyajikan jadwal penelitian dalam bentuk bagan berbasis bulan; durasi dan periode pelaksanaan akan disesuaikan dengan jadwal yang ditetapkan.
